\chapter{第一章 质点运动学}

\section{运动时空测量}
    \subsection{物质与运动}
    \begin{enumerate}
        \item 质点运动学：描述机械运动物体空间位置随时间的变化关系，不探究运动产生与改变的原因。
        \item 物质定义：客观存在，能量的表达；以\textbf{实物}和\textbf{场}两种基本形式存在，运动是物质固有属性，场是更基本的物质形态。
    \end{enumerate}

    \subsection{测量的标准}
    \begin{enumerate}
        \item 物理量测量规则：将物理量与公认标准比对，单位是物理量的命名，标准为精确不变的基准。
        \item 标准选取原则：易于获取、恒定不变。
    \end{enumerate}

    \subsection{国际单位制与量纲}
    \subsubsection{SI国际单位制7个基本量}
    \begin{enumerate}
        \item 长度 $\mathrm{L}$，单位：米 $\mathrm{m}$
        \item 时间 $\mathrm{T}$，单位：秒 $\mathrm{s}$
        \item 质量 $\mathrm{M}$，单位：千克 $\mathrm{kg}$
        \item 电流 $\mathrm{I}$，单位：安培 $\mathrm{A}$
        \item 热力学温度 $\varTheta$，单位：开尔文 $\mathrm{K}$
        \item 物质的量 $\mathrm{N}$，单位：摩尔 $\mathrm{mol}$
        \item 发光强度 $\mathrm{J}$，单位：坎德拉 $\mathrm{cd}$
    \end{enumerate}
    由基本量导出的物理量称为\textbf{导出量}，对应单位为导出单位。

    \subsubsection{量纲}
    物理量用基本量乘方乘积表达的形式；\textbf{量纲一致性原理}：物理方程等式两侧量纲必须完全相同，可用于校验公式正误。
    例：力的量纲 $\mathrm{MLT^{-2}}$；所有量纲指数为 $0$ 的量为无量纲量。

    \subsection{时间与空间}
    \begin{enumerate}
        \item 时空本质：客观存在，是运动物质的存在形式，二者通过物质运动相互关联，无法脱离物质独立存在。
        \item 时间定义：描述物质运动持续性与顺序性，1967年定义：铯-133原子基态两个超精细能级跃迁辐射9192631770个周期的时长为1秒。
        \item 长度定义：描述物质广延性，1983年定义：光在真空中 $1/299792458\ \mathrm{s}$ 内传播的路径长度为1米。
    \end{enumerate}

    \subsection{参考系和坐标系}
        \subsubsection{参考系}
        运动具有绝对性，观测描述具有相对性；被选定作为观测基准的物体称为\textbf{参考系}，所有物理量必须统一至同一参考系计算。

        \subsubsection{坐标系}
        在选定参考系中搭建的定量描述系统，常用：直角坐标系、自然坐标系、极坐标系、球坐标系、柱坐标系；坐标系选定等价于参考系选定，选取原则简化运算。

\section{质点运动的描述}
    \subsection{质点}
    理想化物理模型：当物体自身尺寸、形状对研究问题无影响时，抽象为集中全部质量的几何点。
    相对性示例：地球公转可视为质点，地球自转不可视为质点。

    \subsection{位置矢量与运动方程}
    \subsubsection{位置矢量（位矢）}
    参考系原点指向质点空间位置的矢量，记作 $\boldsymbol{r}$。
    直角坐标系分解：
    \[
    \boldsymbol{r}(t) = x(t)\boldsymbol{i} + y(t)\boldsymbol{j} + z(t)\boldsymbol{k}
    \]
    分量形式：
    \[
    \begin{cases}
    x = x(t)\\
    y = y(t)\\
    z = z(t)
    \end{cases}
    \]
    位矢大小：
    \[
    r = |\boldsymbol{r}| = \sqrt{x^2+y^2+z^2}
    \]
    方向由方向余弦确定：
    \[
    \cos\alpha = \frac{x}{r},\quad \cos\beta = \frac{y}{r},\quad \cos\gamma = \frac{z}{r}
    \]

    \subsubsection{运动方程与轨道方程}
    \begin{enumerate}
        \item 运动方程：$\boldsymbol{r}=\boldsymbol{r}(t)$，描述位矢随时间变化规律；
        \item 轨道方程：消去运动方程参数 $t$ 得到空间轨迹方程，分为直线运动、曲线运动。
    \end{enumerate}

    \subsection{位移与路程}
    \begin{enumerate}
        \item 位移矢量 $\Delta\boldsymbol{r}$：$\Delta t$ 内位矢的增量，矢量，方向由初位置指向末位置
        \[
        \Delta\boldsymbol{r} = \boldsymbol{r}_B - \boldsymbol{r}_A = \Delta x \boldsymbol{i} + \Delta y \boldsymbol{j} + \Delta z \boldsymbol{k}
        \]
        位移大小：
        \[
        |\Delta\boldsymbol{r}| = \sqrt{\Delta x^2+\Delta y^2+\Delta z^2}
        \]
        \item 路程 $\Delta s$：质点实际走过的轨迹长度，标量；
        \item 核心关系：一般 $|\Delta\boldsymbol{r}| \neq \Delta s$，仅单向直线运动或无穷小时间间隔内 $\mathrm{d}s = |\mathrm{d}\boldsymbol{r}|$。
    \end{enumerate}

    \subsection{速度矢量}
    \subsubsection{平均速度}
    $\Delta t$ 内位矢平均变化率，矢量：
    \[
    \bar{\boldsymbol{v}} = \frac{\Delta\boldsymbol{r}}{\Delta t}
    \]

    \subsubsection{瞬时速度（速度）}
    $\Delta t \to 0$ 时平均速度极限，沿轨迹切线方向：
    \[
    \boldsymbol{v} = \lim_{\Delta t \to 0}\frac{\Delta\boldsymbol{r}}{\Delta t} = \frac{\mathrm{d}\boldsymbol{r}}{\mathrm{d}t}
    \]
    直角坐标系分解：
    \[
    \boldsymbol{v} = \frac{\mathrm{d}x}{\mathrm{d}t}\boldsymbol{i} + \frac{\mathrm{d}y}{\mathrm{d}t}\boldsymbol{j} + \frac{\mathrm{d}z}{\mathrm{d}t}\boldsymbol{k} = v_x\boldsymbol{i}+v_y\boldsymbol{j}+v_z\boldsymbol{k}
    \]

    \subsubsection{速率}
    速度矢量的模，标量：
    \[
    v = |\boldsymbol{v}| = \sqrt{v_x^2+v_y^2+v_z^2}
    \]
    注意：平均速度大小 $\neq$ 平均速率。

\section{质点运动变化的描述}
    \subsection{加速度矢量}
    描述速度矢量（大小、方向）随时间变化快慢的物理量。
    \begin{enumerate}
        \item 平均加速度：
        \[
        \bar{\boldsymbol{a}} = \frac{\Delta\boldsymbol{v}}{\Delta t}
        \]
        \item 瞬时加速度：
        \[
        \boldsymbol{a} = \lim_{\Delta t \to 0}\frac{\Delta\boldsymbol{v}}{\Delta t} = \frac{\mathrm{d}\boldsymbol{v}}{\mathrm{d}t} = \frac{\mathrm{d}^2\boldsymbol{r}}{\mathrm{d}t^2}
        \]
    \end{enumerate}
    直角坐标系分解：
    \[
    \boldsymbol{a} = \frac{\mathrm{d}v_x}{\mathrm{d}t}\boldsymbol{i} + \frac{\mathrm{d}v_y}{\mathrm{d}t}\boldsymbol{j} + \frac{\mathrm{d}v_z}{\mathrm{d}t}\boldsymbol{k}
    \]
    加速度大小：
    \[
    a = \sqrt{a_x^2+a_y^2+a_z^2}
    \]

    \subsubsection{两类质点运动学核心问题}
    \begin{enumerate}
        \item 微分问题：已知运动方程 $\boldsymbol{r}(t)$，一阶导求速度，二阶导求加速度；
        \item 积分问题：已知加速度/速度与初始条件，积分逆运算求解速度、运动方程。
    \end{enumerate}
    \subsubsection{运动叠加原理}
    \textbf{运动叠加原理}：任意曲线运动可分解为三个正交方向独立直线运动的叠加。

    \subsection{自然坐标系中的速度和加速度}
    适用于已知轨迹的平面曲线运动，以切向单位矢量 $\boldsymbol{e}_\mathrm{t}$、法向单位矢量 $\boldsymbol{e}_\mathrm{n}$ 搭建坐标系。

    \subsubsection{自然坐标与速度}
    弧长坐标 $s=s(t)$，速率 $v=\dfrac{\mathrm{d}s}{\mathrm{d}t}$，速度矢量：
    \[
    \boldsymbol{v} = \frac{\mathrm{d}s}{\mathrm{d}t}\boldsymbol{e}_\mathrm{t} = v\boldsymbol{e}_\mathrm{t}
    \]

    \subsubsection{切向加速度与法向加速度}
    对速度求导分解：
    \[
    \boldsymbol{a} = \frac{\mathrm{d}v}{\mathrm{d}t}\boldsymbol{e}_\mathrm{t} + \frac{v^2}{\rho}\boldsymbol{e}_\mathrm{n}
    \]
    \begin{enumerate}
        \item 切向加速度：
        \[
        a_\mathrm{t} = \frac{\mathrm{d}v}{\mathrm{d}t}
        \]
        物理意义：改变速度的大小；
        \item 法向加速度（向心加速度）：
        \[
        a_\mathrm{n} = \frac{v^2}{\rho}
        \]
        物理意义：改变速度的方向，$\rho$ 为轨迹曲率半径；
    \end{enumerate}
    合加速度大小：
    \[
    a = \sqrt{a_\mathrm{t}^2 + a_\mathrm{n}^2}
    \]

    \subsubsection{由切向、法向加速度判定运动类型}
    \begin{enumerate}
        \item 若 $a_\mathrm{n} \equiv 0$：质点做直线运动；若 $a_\mathrm{n} \not\equiv 0$：质点做曲线运动；
        \item 若 $a_\mathrm{t} \equiv 0$：质点做匀速率运动；若 $a_\mathrm{t} \not\equiv 0$：质点做变速率运动。
    \end{enumerate}
    直线运动、匀速率曲线运动均为平面一般曲线运动的特例。

    \subsection{圆周运动的角量描述}
    \subsubsection{线量与角量定义}
    \begin{enumerate}
        \item \textbf{线量}：自然坐标系描述物理量，包括弧长 $s$、路程 $\Delta s$、速率 $v$、切向加速度 $a_\mathrm{t}$、法向加速度 $a_\mathrm{n}$；
        \item \textbf{角量}：极坐标系描述圆周运动物理量，包括角位置 $\theta$、角位移 $\Delta\theta$、角速度 $\omega$、角加速度 $\alpha$。
    \end{enumerate}
    规定：质点绕原点逆时针转动角量为正值，顺时针转动角量为负值。

    \subsubsection{角速度与角加速度定义式}
    角速度：
    \[
    \omega = \lim_{\Delta t \to 0}\frac{\Delta\theta}{\Delta t} = \frac{\mathrm{d}\theta}{\mathrm{d}t}
    \]
    角加速度：
    \[
    \alpha = \lim_{\Delta t \to 0}\frac{\Delta\omega}{\Delta t} = \frac{\mathrm{d}\omega}{\mathrm{d}t} = \frac{\mathrm{d}^2\theta}{\mathrm{d}t^2}
    \]
    SI单位：角速度 $\mathrm{rad\cdot s^{-1}}$，角加速度 $\mathrm{rad\cdot s^{-2}}$。

    \subsubsection{圆周运动角量与线量关系式}
    \[
    \begin{cases}
    v = R\omega \\
    a_\mathrm{t} = R\alpha \\
    a_\mathrm{n} = R\omega^2 = \dfrac{v^2}{R}
    \end{cases}
    \]

    \subsubsection{角速度与角加速度的矢量表达}
    无限小角位移满足矢量运算法则，瞬时角速度 $\boldsymbol{\omega}$、瞬时角加速度 $\boldsymbol{\alpha}$ 为矢量，方向由右手螺旋定则判定。
    线速度矢量叉乘公式：
    \[
    \boldsymbol{v} = \boldsymbol{\omega} \times \boldsymbol{r}
    \]
    加速度矢量叉乘公式：
    \[
    \boldsymbol{a} = \boldsymbol{\alpha} \times \boldsymbol{r} + \boldsymbol{\omega} \times \boldsymbol{v}
    \]
    式中第一项为切向加速度，第二项为法向加速度。

\section{相对运动}
    \subsection{参考系平动变换基本关系}
    设参考系 $S'$ 相对惯性系 $S$ 做平动，$S'$ 原点位矢为 $\boldsymbol{r}_0$，质点在 $S$ 系位矢 $\boldsymbol{r}$，在 $S'$ 系位矢 $\boldsymbol{r}'$，满足位矢变换：
    \[
    \boldsymbol{r} = \boldsymbol{r}' + \boldsymbol{r}_0
    \]
    位移变换关系：
    \[
    \Delta\boldsymbol{r} = \Delta\boldsymbol{r}' + \Delta\boldsymbol{r}_0
    \]

    \subsection{伽利略速度变换定理}
    对位矢公式对时间求导，令质点相对 $S$ 速度 $\boldsymbol{v}$，质点相对 $S'$ 速度 $\boldsymbol{v}'$，$S'$ 相对 $S$ 牵连速度 $\boldsymbol{u}$：
    \[
    \boldsymbol{v} = \boldsymbol{v}' + \boldsymbol{u}
    \]
    该式为\textbf{伽利略速度变换公式}，适用于物体运动速度远小于光速的宏观低速场景。

    \subsection{伽利略加速度变换定理}
    对速度变换式再次求导，$S'$ 相对 $S$ 牵连加速度 $\boldsymbol{a}_0$，质点在两参考系加速度 $\boldsymbol{a}$、$\boldsymbol{a}'$：
    \[
    \boldsymbol{a} = \boldsymbol{a}' + \boldsymbol{a}_0
    \]
    若 $S'$ 相对 $S$ 做匀速直线平动，牵连加速度 $\boldsymbol{a}_0 = \boldsymbol{0}$，则：
    \[
    \boldsymbol{a} = \boldsymbol{a}'
    \]
    即惯性系中所有匀速平动参考系内观测到的质点加速度完全相等。